\documentclass[11pt]{article}
\usepackage{amsmath,amssymb,geometry}
\geometry{margin=22mm}
\begin{document}

\section*{P115: completed finite-horizon ridesharing MDP}
The model is a finite-horizon tabular MDP with the transition kernel and reward table supplied by the instance.

\paragraph{Data and domains.}
There are three regions and 180 vehicles in the completed scenario.
$H=12$ hourly decision stages are $t=0,\ldots,11$; thus the original inclusive
terminal reward index is $T=H-1=11$.
The nonterminal states are $s=(t,u,v)$ with $u,v\in\{0,1,2\}$,
the demand and vehicle-availability regimes. Each state's request, vacant,
occupied and previous-demand quantities are explicitly stored in
\texttt{model\_data.states}. Add one absorbing terminal state $\dagger$.
There are $108+1=109$ states.
The action sets contain all regional price/wage combinations
\[
\mathcal A=\{2.4,2.8,3.3\}^{3}\times\{1.3,1.65\}^{3},
\qquad |\mathcal A|=216.
\]
The action list order is the index used by each state's transition table.
Let $r_{sa}$ be the table's expected net revenue, and let $P_{sas'}$
be its next-state probability. The table also records expected served counts
for reporting.
The reward table is the completed income-minus-wage reward, with finite
availability and heterogeneous demand/trip inputs recorded in its generator.
$P_{sas'}\ge0$ and $\sum_{s'}P_{sas'}=1$.
Every nonterminal transition advances the clock by one, and all transitions
from $t=11$ end at $\dagger$. All these entries are fixed input data.
Initial distribution $\mu$ is explicitly supplied on the nine $t=0$ states,
$\mu\ge0$, $\sum_s\mu_s=1$; $\gamma=0.97$.

\paragraph{Policy optimization.}
Choose one distribution $\pi(\cdot\mid s)$ for each clock-augmented state:
\[
\pi(a\mid s)\ge0,\qquad \sum_{a\in\mathcal A}\pi(a\mid s)=1.
\]
With $V^\pi_\dagger=0$, impose all Bellman equalities
\[
V^\pi_s=\sum_a\pi(a\mid s)
 \left(r_{sa}+\gamma\sum_{s'}P_{sas'}V^\pi_{s'}\right)
 \quad(s\ne\dagger),
\qquad
\max_\pi J=\sum_s\mu_sV^\pi_s.
\]
Equivalently,
$J=E_\pi[\sum_{t=0}^{11}\gamma^t r_{s_ta_t}]$.
Clock augmentation makes a stationary table $\pi(a\mid s)$ capable of
representing every admissible nonstationary policy in the original
finite-horizon state description.

\paragraph{Exact solution and verification.}
Reverse-time dynamic programming computes
\[
V^*_\dagger=0,\qquad
V^*_s=\max_a\{r_{sa}+\gamma\sum_{s'}P_{sas'}V^*_{s'}\},
\qquad \pi^*_s\in\arg\max_a(\cdot).
\]
A deterministic maximizing action exists because each action set is finite.
The implementation separately propagates $\mu$ forward under that policy,
recomputes the discounted objective and tests every alternative action's
Bellman inequality. Transition normalization and terminal probability are
also checked. The recorded value is an exact finite-table optimum up to
floating-point roundoff for the stated finite tabular model.

\end{document}
